\begin{table}[htbp]
\centering
\caption{Summary Statistics}
\label{tab:summary}
\begin{adjustbox}{max width=\textwidth}
\begin{tabular}{lcccc}
\hline\hline
 & Mean & SD & Median & N \\
\hline
\multicolumn{5}{l}{\textit{Panel A: Full Sample}} \\
Vote margin (\%) & 12.77 & 41.8 & 5.79 & 6268 \\
Wealthier candidate won & 0.597 & 0.491 &  & 6268 \\
Rich candidate assets (Rs lakhs) & 436.8 & 3917.2 & 99.8 & 6268 \\
Poor candidate assets (Rs lakhs) & 65.4 & 210.2 & 20.7 & 6268 \\
Wealth ratio & 90.8 & 1584.2 & 3.9 & 6268 \\
Log wealth difference & 1.77 & 1.52 & 1.36 & 6268 \\
Rich candidate criminal cases & 0.61 & 1.7 & 0 & 6268 \\
Poor candidate criminal cases & 0.6 & 1.85 & 0 & 6268 \\
Reserved constituency & 0.264 & 0.441 &  & 6268 \\
Number of states & 29 &  &  &  \\
Number of elections & 10 &  &  &  \\
\hline
\multicolumn{5}{l}{\textit{Panel B: Close Elections ($|$margin$|$ $<$ 5\%)}} \\
Vote margin (\%) & $-$0.14 & 2.9 & $-$0.15 & 1082 \\
Wealthier candidate won & 0.481 & 0.500 &  & 1082 \\
Rich candidate assets (Rs lakhs) & 310.7 & 944.8 & 98.0 & 1082 \\
Poor candidate assets (Rs lakhs) & 62.6 & 120.5 & 25.8 & 1082 \\
Wealth ratio & 29.3 & 399.2 & 3.2 & 1082 \\
\hline\hline
\end{tabular}
\end{adjustbox}
\begin{minipage}{0.95\textwidth}
\footnotesize
\textit{Notes:} Summary statistics for the RDD analysis sample of constituency-elections where both the top-two vote-getters have affidavit data. The ``wealthier candidate'' is the one with higher total declared assets. Vote margin is defined as the percentage point difference between the wealthier and poorer candidate's vote shares. Panel B restricts to elections where the absolute vote margin is within 5 percentage points. Wealth ratio = rich candidate's assets / poor candidate's assets.
\end{minipage}
\end{table}
